\section{总结与延伸}

Week 7 探讨了代码审查在 AI 编码时代的角色演变。核心信息是：AI 让代码审查更高效，但同时也让代码审查更加不可或缺——因为你需要审查的不仅是人写的代码，还有 AI 写的代码。

\subsection{关键要点}

\begin{itemize}
    \item 代码审查的错误检出率（55--60\%）远超测试（25--45\%）
    \item 好的审查意见是建设性的、具体的，并邀请对话
    \item 审查应覆盖五个维度：正确性、可读性、性能、安全、最佳实践
    \item AI 代码审查适合处理``表层''问题，人类专注``深层''问题
    \item 当前局限：假阳性、无法捕获 repo-specific idioms、无法处理复杂业务逻辑
    \item AI 时代代码审查更重要，``AI wrote it'' 不是免责借口
\end{itemize}

\subsection{团队落地建议}

\begin{enumerate}
    \item 先把 PR 模板、变更说明和测试说明标准化，再引入 AI review
    \item 让 AI 先覆盖高频表层问题，不要一开始就要求它替代资深 reviewer
    \item 为关键评论建立证据链：引用 diff、规则或历史上下文
    \item 持续统计假阳性与漏报率，把 review 系统当产品运营
    \item 在安全、支付、权限等高风险模块保留更强的人类审批门槛
\end{enumerate}

\subsection{审查策略对照表}

\begin{table}[H]
\centering
\begin{tabular}{p{0.24\textwidth} p{0.27\textwidth} p{0.31\textwidth}}
\toprule
场景 & 更适合 AI 的任务 & 更适合人类的任务 \\
\midrule
常规业务 PR & 命名、重复逻辑、常见 bug、规则检查 & 业务意图、边界取舍、代码演进方向 \\
基础设施改造 & 配置校验、依赖影响扫描 & 回滚策略、兼容性风险、迁移节奏 \\
安全敏感模块 & 已知漏洞模式匹配、权限变更提示 & 威胁建模、攻击面判断、例外审批 \\
大规模重构 & 模式一致性检查、遗漏引用发现 & 是否值得重构、拆分策略、上线窗口 \\
\bottomrule
\end{tabular}
\caption{AI reviewer 与人类 reviewer 的分工边界}
\end{table}

\subsection{拓展阅读}

\begin{itemize}
    \item Graphite：\url{https://graphite.dev}
    \item Greptile：\url{https://greptile.com}
    \item CodeRabbit：\url{https://coderabbit.ai}
    \item Google 的代码审查最佳实践：\url{https://google.github.io/eng-practices/review/}
    \item ``How to do a code review'' by Blake Smith
\end{itemize}

%% --- Fin del contenido --- %%

\end{document}
